\title{xLSTM: Extended Long Short-Term Memory}

\begin{abstract}
The foundational concepts of the constant error carousel and gating, introduced in the 1990s, underpin the architecture of Long Short-Term Memory (LSTM) networks. These mechanisms have established LSTMs as enduringly influential in the landscape of deep learning, notably serving as the foundation for the earliest Large Language Models (LLMs). Nonetheless, with the emergence of the Transformer architecture—distinguished by its scalable, parallelizable self-attention mechanism—a significant paradigm shift occurred, leading to Transformers surpassing LSTMs at larger scales. This work revisits a fundamental inquiry: What level of performance can be attained in language modeling if LSTMs are scaled to billions of parameters and enhanced by state-of-the-art LLM techniques while also addressing their well-known shortcomings? To this end, we propose exponential gating mechanisms supported by tailored normalization and stabilization strategies. We also revise the structure of LSTM memory, introducing: (i) the sLSTM variant, featuring scalar memory, scalar updates, and novel schemes for memory integration, and (ii) the mLSTM, a fully parallelizable design that utilizes a matrix-based memory and an update rule grounded in covariance. By embedding these enhanced LSTM designs into residual block frameworks, we create xLSTM modules, which are then sequentially layered to form full xLSTM networks. The adoption of exponential gating and redesigned memory representations enables xLSTM to compare favorably with leading Transformer models and State Space Models, achieving competitive or superior results in both effectiveness and scalability.\\
Code available at: \url{https://github.com/NX-AI/xlstm}
\end{abstract}

\section{Introduction}
\label{sec:intro}

The Long Short-Term Memory (LSTM) framework~\citep{Hochreiter:91,Hochreiter:97nips,Hochreiter:97}, characterized by features such as the constant error carousel and gating mechanisms, was developed as a solution to the vanishing gradient issue present in recurrent neural networks~\citep{Hochreiter:91,Hochreiter:00book}:

\begin{align}
\label{eq:lstm_idea}
  c_t \ = \ f_t \ c_{t-1} \ + \ 
          i_t \  z_t \ , \quad  
          h_t \ = \ o_t \ \psi({c_t}) \ . 
\end{align}

The constant error carousel involves an additive modification to the cell state $c_{t-1}$ (depicted in green), where new inputs $z_t$ are incorporated, and this process is regulated by sigmoid-based gates (shown in blue). Specifically, the input gate $i_t$ and the forget gate $f_t$ govern how the cell state is updated, whereas the output gate $o_t$ determines the resultant output from the memory cell—namely, the hidden state $h_t$. Before producing the hidden state, the cell state undergoes normalization or is squashed using $\psi$, after which the output gate applies further gating to yield the final hidden state.

LSTMs have achieved notable success across a broad range of domains \citep{Hochreiter:01,Hochreiter:07,Schmidhuber:15}, dominating the landscape of text generation prior to the emergence of Transformers in 2017~\citep{Vaswani:17}. Their utility has been established in multiple sequence-oriented applications, including but not limited to text generation~\citep{Graves:13,Karpathy:15blog}, handwriting synthesis~\citep{Graves:13}, sequence-to-sequence translation tasks~\citep{Sutskever:14nips}, program evaluation~\citep{Zaremba:14arxiv}, automatic image captioning~\citep{Karpathy:15,Hossain:19}, code generation~\citep{Karpathy:15blog}, modeling rainfall-runoff processes~\citep{Kratzert:18,Kratzert:19}, and in hydrological modeling for flood prediction~\citep{Nearing:24}. Within reinforcement learning, LSTMs have outperformed alternative sequence models in high-profile applications, including the AlphaStar agent for StarCraft II~\citep{Vinyals:17short}, the OpenAI Five system for Dota 2~\citep{OpenAI:19}, and controllers used in nuclear fusion experiments~\citep{Degrave:22short}. The strength of LSTMs lies in their capacity for abstraction, enabling them to proficiently capture and retain semantic information within their memory units~\citep{Karpathy:15blog}, as demonstrated by the identification of neurons corresponding to numeracy and syntax~\citep{Lakretz:19}, language properties~\citep{Bau:19}, and sentiment~\citep{Radford:17arxiv}. Despite advancements in model architectures, LSTMs continue to be employed in mission-critical scenarios~\citep{Degrave:22short,Nearing:24} and have maintained their relevance over time.

Although LSTMs have achieved significant advances, they are hindered by three primary drawbacks: (i) Lack of ability to update previous storage choices. This issue can be illustrated through the {\it Nearest Neighbor Search} task (see also Appendix~\ref{sec:appNearestNeighborSearch}): Given a reference vector, the sequence is processed in order to identify and retrieve the value associated with the closest vector at the end of the sequence. The left panel in Figure~\ref{fig:lstmProblems} presents the average squared error for this task. Standard LSTM architectures encounter difficulty overwriting a memorized value upon discovering a more similar vector, whereas the proposed xLSTM overcomes this with exponential gating. (ii) Constrained memory capacity, as all data must be encoded in scalar cell states. This is demonstrated by the {\it Rare Token Prediction} problem. The right panel of Figure~\ref{fig:lstmProblems} shows token prediction perplexity on Wikitext-103~\citep{Merity:17}, split by token frequency. The LSTM's limited capacity leads to inferior performance on infrequent tokens. The xLSTM mitigates this constraint by utilizing a matrix-based memory structure. (iii) Poor suitability for parallel computation, a consequence of memory mixing from recurrent hidden-to-hidden connections between subsequent time steps, which necessitates strict sequential processing.

The shortcomings inherent to LSTM architectures have facilitated the rise of Transformers~\citep{Vaswani:17} within the field of language modeling. An open question remains: if these LSTM constraints are addressed and LSTMs are scaled comparably to modern Large Language Models, what levels of performance might they attain in language modeling tasks?

\section{Extended Long Short-Term Memory}
\label{sec:xLSTM}

To address the constraints inherent in LSTM, Extended Long Short-Term Memory (xLSTM) introduces two principal innovations to the foundational LSTM framework outlined in Equation~\eqref{eq:lstm_idea}. These enhancements—exponential gating mechanisms and innovative memory architectures—expand the LSTM paradigm by introducing two novel architectures: (i) sLSTM (refer to Section~\ref{sec:sLSTM}), which employs a scalar-valued memory, scalar updates, and incorporates memory mixing; and (ii) mLSTM (detailed in Section~\ref{sec:mLSTM}), which utilizes a matrix-form memory along with updates based on a covariance (outer product) formulation, facilitating full parallel computation. Both of these new architectures—sLSTM and mLSTM—integrate exponential gating for improved performance. Notably, in pursuit of parallelism, mLSTM dispenses with memory mixing (i.e., removes hidden-hidden recurrence). The design of both mLSTM and sLSTM can be naturally generalized to encompass multiple memory cells, with sLSTM supporting memory mixing between cells. Additionally, sLSTM may adopt a multi-head structure wherein memory mixing occurs only among cells within a given head and not across different heads, introducing a novel approach to memory mixing when paired with exponential gating. Conversely, for mLSTM, employing multiple heads is functionally indistinguishable from using multiple memory cells.

By embedding these novel LSTM variants within residual block structures, we construct what we term xLSTM blocks (refer to Section~\ref{sec:xLSTMblock}). Arranging these xLSTM blocks in a residual stacking manner leads to the formation of xLSTM architectures, as detailed in Section~\ref{sec:xLSTMarchitecure}. The components and layout of the xLSTM architecture are illustrated in Figure~\ref{fig:xlstm_sketch}.

\subsection{Review of the Long Short-Term Memory}
\label{sec:orgLSTM}

The foundational concept of LSTM~\citep{Hochreiter:91,Hochreiter:97nips,Hochreiter:97} established the scalar memory cell as the core unit for both computation and information retention, effectively mitigating the problem of vanishing gradients~\citep{Hochreiter:91,Hochreiter:00book} via the constant error carousel, or cell state update mechanism. Within the memory cell reside three distinct gates: the input gate, output gate, and forget gate. The introduction of the forget gate is credited to \citet{Gers:00}.  
The evolution of the LSTM memory cell across time step $t$ can be described by the following update equations:

\begin{align}
c_t \ &= \  f_t \ c_{t-1} \ + \ i_t \ z_t &  & &\text{cell state} \\
h_t  \ &= \ o_t \ \tilde{h}_t \ , & \tilde{h}_t \ &= \ \psi \left( c_t\right)
  &\text{hidden state} \\
z_t \ &= \ \varphi \left( \tilde{z}_t \right) \ , 
  &\tilde{z}_t \ &=  \ \mathbf{w}^\top_{z} \ \mathbf{x}_t \ + \
  r_{z}  h_{t-1} \ + \  b_{z} \ \
  &\text{cell input} \\
i_t \ &= \ \sigma \left( \tilde{i}_t  \right) \ , 
  &\tilde{i}_t \ &= \ \mathbf{w}^\top_{i} \ \mathbf{x}_t \ + \
  r_{i}  \ h_{t-1} \ + \  b_{i} \ \
  &\text{input gate} \\
f_t \ &= \ \sigma \left(  \tilde{f}_t \right) \ , 
  &\tilde{f}_t \ &= \ \mathbf{w}^\top_{f} \ \mathbf{x}_t  \ + \
  r_{f}  \ h_{t-1} \ + \  b_{f} \ \
  &\text{forget gate} \\
o_t \ &= \ \sigma \left( \tilde{o}_t \right) \ , 
  &\tilde{o}_t  \ &= \ \mathbf{w}^\top_{o} \ \mathbf{x}_t \ + \
  r_{o}  \ h_{t-1} \ + \  b_{o} \ \
  &\text{output gate} 
\end{align}

The vectors $\mathbf{w}_{z}$, $\mathbf{w}_{i}$, $\mathbf{w}_{f}$, and $\mathbf{w}_{o}$ represent the weights connecting the input $\mathbf{x}_t$ to the cell input, input gate, forget gate, and output gate, in that order. The parameters $r_{z}$, $r_{i}$, $r_{f}$, and $r_{o}$ serve as the recurrent weight coefficients linking the previous hidden state $h_{t-1}$ to the cell input, input gate, forget gate, and output gate, respectively. The terms $b_{z}$, $b_{i}$, $b_{f}$, and $b_{o}$ denote the corresponding biases. The activation functions for the cell input and hidden state are represented by $\varphi$ and $\psi$, commonly instantiated as $\tanh$. Notably, $\psi$ is applied to ensure the cell state values are confined within a finite range. All gate activations employ the sigmoid nonlinearity, i.e., $\sigma \left( x \right)= 1/(1+\exp(-x))$. Subsequent developments in the architecture introduced a vectorized form of memory cells, $\mathbf{c}_t \in \mathbb{R}^{d}$, which facilitates the application of recurrent weight matrices $\mathbf{R} \in \mathbb{R}^{d \times d}$ to interactively mix the memory cell outputs~\citep{Greff:15}; refer to Appendix~\ref{sec:apporgLSTM} for further elaboration. Empirical ablation analyses have demonstrated the essential nature of all memory cell elements~\citep{Greff:15}.

\subsection{sLSTM}
\label{sec:sLSTM}

To enable LSTMs to modify their storage choices, we incorporate exponential gates (depicted in red) along with mechanisms for normalization and stabilization. Specifically, both the input and forget gates are permitted to employ exponential activation functions. For the normalization process, we define a normalizer state that accumulates the sum of the input gate values multiplied by all subsequent forget gates.

The forward pass for scalar sLSTM is given by:

\begin{align}
\label{eq:slstmforward}
c_t \ &= \  f_t \ c_{t-1} \ + \ i_t \ z_t & & 
 &\text{cell state} \\
n_t \ &= \  f_t \ n_{t-1} \ + \ i_t  & & 
 &\text{normalizer state} \\
h_t \ &= \ o_t \ \tilde{h}_t \ ,
  &\tilde{h}_t \ &= \ c_t / n_t
&\text{hidden state} \\
z_t \ &= \ \varphi \left( \tilde{z}_t \right) \ , 
  &\tilde{z}_t \ &=  \ \mathbf{w}^\top_{z} \ \mathbf{x}_t \ + \
  r_{z}  h_{t-1} \ + \  b_{z}
  &\text{cell input} \\
i_t \ &= \ \!\exp \left( \tilde{i}_t  \right) \ , 
  &\tilde{i}_t \ &= \ \mathbf{w}^\top_{i} \ \mathbf{x}_t \ + \
  r_{i}  \ h_{t-1} \ + \  b_{i}
  &\text{input gate} \\
f_t \ &= \ \sigma \left(  \tilde{f}_t \right) \ \text{OR} \
\!\exp \left(  \tilde{f}_t \right) \ ,
  &\tilde{f}_t \ &= \ \mathbf{w}^\top_{f} \ \mathbf{x}_t  \ + \
  r_{f}  \ h_{t-1} \ + \  b_{f}
  &\text{forget gate} \\
o_t \ &= \ \sigma \left( \tilde{o}_t \right) \ , 
  &\tilde{o}_t  \ &= \ \mathbf{w}^\top_{o} \ \mathbf{x}_t \ + \
  r_{o}  \ h_{t-1} \ + \  b_{o}
  &\text{output gate} 
\end{align}

We adapt the gating mechanisms from the standard LSTM framework—including gates that depend on inputs or hidden states, as well as the associated bias terms—to the novel architectures. Since exponential activation functions may produce excessively large outputs resulting in numerical overflows, we enhance gate stability by incorporating an extra state denoted by \mbox{$m_t$~\citep{Milakov:18arxiv}}:

\begin{align}
\label{eq:slstmstabil}
m_t \ &= \ \max \left( \log ( f_t ) + m_{t-1} , \log ( i_t ) \right) &\text{stabilizer state} \\
i'_t \ &= \ \!\exp \left( \log \left ( i_t \right) - m_t \right) \ = \ \!\exp \left( \tilde{i}_t - m_t \right) \ \, 
  &\text{stabil. input gate} \\
f'_t \ &= \ \!\exp \left( \log \left( f_t \right) + m_{t-1} - m_t \right) \ \, 
  &\text{stabil. forget gate}
\end{align}

Appendix~\ref{sec:appsLSTM} demonstrates that substituting $f_t$ with $f'_t$ and $i_t$ with $i'_t$ during the forward computation leaves both the network’s output and the gradients of the loss with respect to the parameters unchanged.

\paragraph{New Memory Mixing.} Similar to the standard LSTM, sLSTM supports multiple memory cells (refer to Appendix~\ref{sec:appsLSTM}). The use of multiple memory cells facilitates mixing of memory via recurrent pathways $\mathbf{R}_{\mathbf{z}}$, $\mathbf{R}_{\mathbf{i}}$, $\mathbf{R}_{\mathbf{f}}$, $\mathbf{R}_{\mathbf{o}}$ originating from the hidden state vector $\mathbf{h}$ and connecting to both the memory cell input $\mathbf{z}$ and the gates $\mathbf{i}$, $\mathbf{f}$, $\mathbf{o}$. One distinctive feature introduced in this type of memory mixing is the influence of exponential gating. The updated sLSTM framework permits several heads, where memory mixing is allowed within a head, but is not permitted between different heads. By combining the use of heads in sLSTM with exponential gating, a novel approach to memory mixing is established.

\subsection{mLSTM}
\label{sec:mLSTM}

In order to improve the storage capabilities of LSTM networks, we replace the traditional scalar memory cell $c \in \mathbb{R}$ with a matrix-form memory $\mathbf{C} \in \mathbb{R}^{d \times d}$. This modification enables retrieval operations to be conducted through matrix multiplication. Specifically, at each time step $t$, our goal is to encode a pair of $d$-dimensional vectors, the key $\mathbf{k}_t \in \mathbb{R}^d$ and the value $\mathbf{v}_t \in \mathbb{R}^d$, adopting nomenclature from the Transformer architecture. Subsequently, at a future step $t + \tau$, the intention is to recover the value $\mathbf{v}_t$ using a corresponding query vector $\mathbf{q}_{t+\tau} \in \mathbb{R}^d$. This framework aligns with the principles of Bidirectional Associative Memories (BAMs) \citep{Kohonen:72,Anderson:72,Nakano:72,Anderson:77}.

To encode a key-value association in memory, we utilize the covariance update mechanism~\citep{Sejnowski:77,Dayan:91}, which is given by

\begin{align}
\mathbf{C}_t \ &= \ \mathbf{C}_{t-1} \ + \ \mathbf{v}_{t} \ \mathbf{k}_t^\top \ .
\end{align}

We consider the application of layer normalization prior to mapping the input to key and value representations, ensuring these vectors have zero mean. The covariance update mechanism achieves optimal performance~\citep{Dayan:91} in maximizing the distinguishability of recovered binary vectors, corresponding to an optimal signal-to-noise ratio. This degree of separability can be further increased by constraining retrieval operations to pairwise interactions, albeit at the cost of quadratic computational complexity similar to attention mechanisms~\citep{Krotov:16,Krotov:17,Ramsauer:21}. The covariance update procedure matches that used in Fast Weight Programmers \citep{Schmidhuber:92ncfastweights,Schlag:21}, which were later modified by introducing a fixed decay coefficient for $\mathbf{C}_{t-1}$ and a constant learning rate for 
$\mathbf{v}_{t} \mathbf{k}_t ^\top$~\citep{Ba:16}.
Inspired by these approaches, we incorporate the covariance update rule within the LSTM architecture, interpreting the forget gate as implementing the decay rate, the input gate as setting the learning rate, and the output gate as modulating the retrieved representation.

In this matrix memory approach, the normalizer state is calculated as a weighted combination of key vectors, with each key vector's contribution determined by the input gate and the product of all subsequent forget gates. As before, the normalizer state tracks the accumulated influence of the gates. Because the dot product between the query and the normalizer state may approach zero, we take its absolute value and enforce a minimum threshold (usually set to 1.0), following the strategy from previous work~\citep{Sun:23arxiv}. The forward computation for mLSTM is:

\begin{align}
\label{eq:mlstm_recurrent_begin}
\mathbf{C}_t \ &= \  f_t \ \mathbf{C}_{t-1} \ + \ 
  i_t \ \mathbf{v}_t \ \mathbf{k}_t^\top & &  &\text{cell state} \\
\mathbf{n}_t \ &= \  f_t \ \mathbf{n}_{t-1} \ + \ 
  i_t \ \mathbf{k}_t & &  &\text{normalizer state} \\
\mathbf{h}_t  \ &= \ \mathbf{o}_t \ \odot \ \tilde{\mathbf{h}}_t
\ , \qquad \qquad 
\tilde{\mathbf{h}}_t \ = \   \mathbf{C}_t \mathbf{q}_t \ / \ 
  \max \left\{ |\mathbf{n}_t^\top \mathbf{q}_t|, 1 \right\} 
   &&&\text{hidden state} \\
\mathbf{q}_t \ &= \ \mathbf{W}_q \ \mathbf{x}_t \ + \ \mathbf{b}_q  & &  &\text{query input} \\
\mathbf{k}_t \ &= \ \frac{1}{\sqrt{d}} \mathbf{W}_k \ \mathbf{x}_t \ + \ \mathbf{b}_k  & &  &\text{key input} \\
\mathbf{v}_t \ &= \ \mathbf{W}_v \ \mathbf{x}_t \ + \ \mathbf{b}_v  & &  &\text{value input} \\
i_t \ &= \ \!\exp \!  \! \left( \tilde{i}_t  \right) \ , 
 \qquad \qquad \ \ \ \ \,
  \tilde{i}_t \ = \ \mathbf{w}^\top_{i} \ \mathbf{x}_t \ + \  b_{i} \ \ \ 
  &&&\text{input gate} \\
f_t \ &= \ \sigma \!  \left(  \tilde{f}_t \right) \ \text{OR} \ \!\exp \!  \! \left(  \tilde{f}_t \right) , \ \ \: \!
\tilde{f}_t \ = \ \mathbf{w}^\top_{f} \ \mathbf{x}_t  \ + \
  b_{f}
  &&& \text{forget gate} \\
\label{eq:mlstm_recurrent_end}
\mathbf{o}_t \ &= \ \sigma \left( \tilde{\mathbf{o}}_t \right) \ , \qquad \qquad \quad \ \ \ \,
  \tilde{\mathbf{o}}_t  \ = \ \mathbf{W}_{\mathbf{o}} \ \mathbf{x}_t \ + \
  \mathbf{b}_{\mathbf{o}}
  &&&\text{output gate} 
\end{align}

mLSTM is capable of utilizing several memory cells, similar to the structure of the original LSTM. In the case of mLSTM, employing multiple heads is functionally the same as using multiple cells because memory states remain independent without any interaction. To ensure the stability of the mLSTM’s exponential gates, we adopt the stabilization procedures applied to sLSTM (refer to Equation~\ref{eq:slstmstabil}). 
Given that mLSTM does not incorporate memory mixing, its recurrence relation can be reformulated in a parallelizable manner. Further information can be found in Appendix~\ref{sec:appmLSTM}.

\subsection{xLSTM Architecture}
\label{sec:xLSTMarch}

\paragraph{xLSTM Blocks.}
\label{sec:xLSTMblock}
The role of an xLSTM block is to produce a non-linear, high-dimensional representation of past information, enhancing the distinction between diverse histories or contexts. This differentiation of prior contexts is fundamental for accurately forecasting the subsequent element in a sequence, such as the next token. Our approach leverages Cover’s Theorem~\citep{Cover:65}, which posits that patterns, when embedded non-linearly in a space of higher dimensionality, become more amenable to linear separation compared to the original lower-dimensional setting. We examine two types of residual block configurations: (i) The post up-projection residual block, inspired by the Transformer architecture, first generates a non-linear summary of the past within the initial space, then applies a linear projection into a space of greater dimensionality, follows with a non-linear activation, and subsequently maps back linearly to the starting dimensionality; this setup is illustrated in the left portion of Figure~\ref{fig:Backbones} as well as in the third column of Figure~\ref{fig:xlstm_sketch}. For an expanded illustration, refer to Figure~\ref{fig:appsLSTM_detailed} in the appendix. (ii) The pre up-projection residual block, which resembles approaches found in State Space Models, performs an initial linear projection into a high-dimensional space,

The model first applies a non-linear compression of past information within a high-dimensional space, followed by a linear transformation that projects this representation back to the original space. In the architecture, when an xLSTM block incorporates an sLSTM, the up-projection occurs after this process, using the post up-projection block. Conversely, for an xLSTM block with an mLSTM, the up-projection is executed before, via the pre up-projection block, to accommodate the increased memory capacity offered by the high-dimensional space. For a visual explanation, please consult the left panel of Figure~\ref{fig:Backbones}, the third column of Figure~\ref{fig:xlstm_sketch}, or refer to Figure~\ref{fig:appsLSTM_detailed} located in the appendix.

\paragraph{xLSTM Architecture. }
\label{sec:xLSTMarchitecure}
The xLSTM model is formed by sequentially stacking basic units in a residual manner~\citep{Srivastava:15,He:16}. Our approach employs the prevalent pre-LayerNorm~\citep{Ba:16b} residual backbone, which is standard in current Large Language Models. The final column of Figure~\ref{fig:xlstm_sketch} illustrates this setup.

\subsection{Memory and Speed Considerations}

Unlike Transformers, xLSTM networks exhibit linear computational complexity and maintain constant memory requirements regardless of sequence length. Due to the compressive nature of xLSTM memory, these networks are particularly advantageous for industrial use cases and edge device deployments.

The mLSTM’s memory mechanism is parameter-free, yet it incurs significant computational overhead due to its reliance on a $d \times d$ matrix memory and updates of the same dimensionality. This design balances increased memory capacity with greater computational demand. However, because these operations are highly parallelizable on GPUs, their impact on actual runtime is minimal.

Although mLSTM supports parallelization similar to FlashAttention~\citep{Dao:22,Dao:23} and GLA~\citep{Yang:23arxiv}, sLSTM lacks parallelizability because of memory mixing resulting from hidden-to-hidden connections. Nonetheless, we introduced an efficient CUDA implementation, optimizing GPU memory down to the register level, which results in sLSTM execution being typically no more than twice as slow as mLSTM.

\section{Related Work}
\label{sec:related}

\paragraph{Linear Attention.} Multiple approaches have been proposed to address the issue of quadratic computational cost in Transformers with respect to context length, aiming to achieve attention mechanisms that scale linearly. The Synthesizer introduces the concept of learning synthetic attention weights, entirely bypassing explicit token--token dependencies~\citep{Tay:20arxiv}. Linformer models self-attention through the use of a low-rank matrix representation, facilitating a linear approximation~\citep{Wang:20arxiv}. The Linear Transformer transforms the core attention process itself into a linear operation \citep{Katharopoulos:20}. Performer employs positive orthogonal random features to produce a linear approximation of the softmax attention~\citep{Choromanski:21}. Alternative strategies, such as those found in Structured Global Convolution (SGConv)~\citep{Li:22} and the Hyena Hierarchy~\citep{Poli:23}, opt to replace attention entirely with efficient, long-range convolution operations.

\paragraph{State Space Models.} In recent years, State Space Models (SSMs) have gained significant traction due to their linear scalability with respect to context length and competitive results relative to Transformer architectures. Among the earliest models introduced was the Structured State Space sequence model (S4)~\citep{Gu:21}, which was subsequently followed by architectures such as the Diagonal State Space (DSS) model~\citep{Gupta:22}, Gated State Space (GSS) models~\citep{Mehta:22}, the S5 model~\citep{Smith:22}, Bidirectional Gated SSM (BiGS)~\citep{Wang:22}, H3 model~\citep{Fu:23}, and Mamba~\citep{Gu:24arxiv}.

\paragraph{Recurrent Neural Networks.} Recurrent Neural Networks (RNNs) have emerged as alternatives to Transformers and attention mechanisms, primarily because their computational complexity scales linearly with sequence length. Deep Linear Recurrent Units (LRUs) incorporated within RNNs have delivered encouraging outcomes for language modeling tasks~\citep{Orvieto:23, De:24arxiv}. Additionally, models such as Hierarchically Gated Linear RNN (HGRN)~\citep{Qin:23} and its successor HGRN2~\citep{Qin:24arxiv} have also demonstrated strong performance. Among the RNN-based methods, the RWKV architecture~\citep{Peng:23arxivshort,Peng:24arxivshort} has established itself as a notable contender, achieving results that are comparable to Transformer-based models.

\paragraph{Gating.}
Gating, a fundamental concept introduced by LSTM, has been revisited and given new interpretations in various contemporary methods. This mechanism appears in numerous recent models, including HGRN~\citep{Qin:23}, HGRN2~\citep{Qin:24arxiv}, Gated Linear Attention (GLA)~\citep{Yang:23arxiv}, Gated State Space (GSS) models~\citep{Mehta:22}, Bidirectional Gated SSM (BiGS)~\citep{Wang:22}, Moving Average Equipped Gated Attention (MEGA)~\citep{Ma:22}, RWKV~\citep{Peng:23arxivshort}, as well as Mamba~\citep{Gu:24arxiv}.

\paragraph{Covariance Update Rule.} In order to increase storage capability, we integrated a matrix memory into the mLSTM cell that operates via a covariance update rule. Several alternative approaches also utilize similar update principles, including Fast Weight Programmers \citep{Schmidhuber:92ncfastweights,Schlag:21}, RWKV-5 and RWKV-6~\citep{Peng:24arxivshort}, Retention~\citep{Sun:23arxiv}, Linear Transformer~\citep{Katharopoulos:20}, and HGRN2~\citep{Qin:24arxiv}.

\paragraph{Most Related.} The models most similar in concept to xLSTM are Retention~\citep{Sun:23arxiv}, RWKV~\citep{Peng:23arxivshort,Peng:24arxivshort}, and HGRN2~\citep{Qin:24arxiv}, as they employ mechanisms such as matrix memory and/or gating. Nevertheless, unlike the recently introduced sLSTM, these models lack support for memory mixing. The inclusion of memory mixing is crucial for addressing state tracking tasks, rendering LSTMs more powerful in this regard than either State Space Models (SSMs) or Transformers~\citep{Merrill:24,Deletang:23}. Accurate state tracking plays a vital role in executing tasks like code evaluation or following entities throughout lengthy narratives.

\paragraph{Residually Stacking Architectures.} As is the case with virtually all modern large-scale deep learning models, xLSTM architectures are formed by stacking modular components with residual connections~\citep{Srivastava:15,He:16}.

This design principle paved the way for the success of deep convolutional networks~\citep{He:16} as well as the rise of Transformers~\citep{Vaswani:17}. The Transformer architecture, in particular, serves as the foundational mechanism behind a wide array of Large Language Models (LLMs), including GPT-3~\citep{Brown:20short}, ChatGPT~\citep{Schulman:22short}, GPT-4~\citep{Achiam:23gpt4short}, Megatron-LM~\citep{Shoeybi:19}, Gopher~\citep{Rae:21short}, ERNIE 3.0 Titan~\citep{Wang:21arxivshort}, GLaM~\citep{Du:21glamshort}, Chinese M6~\citep{Lin:21}, AlexaTM 20B (multilingual)~\citep{Soltan:22}, OPT~\citep{Zhang:22opt}, Chinchilla~\citep{Hoffmann:22short}, BLOOM~\citep{Scao:22arxivshort}, GLM-130B~\citep{Zeng:22arxivshort}, LaMDA~\citep{Thoppilan:22short}, PaLM~\citep{Chowdhery:22short}, Llama~\citep{Touvron:23arxiv}, and Gemini~\citep{GeminiTeam:23arxiv,Reid:24arxiv}.

\section{Experiments}
\label{sec:Exp}

This section presents a comprehensive experimental study of xLSTM, focusing on its performance in language modeling relative to existing approaches. Section~\ref{sec:ExpSyntethic} explores the unique features of xLSTM through controlled synthetic experiments. In Section~\ref{sec:ExpComparison}, we analyze and compare the validation set perplexities obtained by a range of state-of-the-art language modeling techniques, each trained on a 15B-token subset of the SlimPajama corpus~\citep{Soboleva:23}. We also carry out ablation studies for xLSTM using the same dataset. Following this, we examine the scaling properties of these models, taking cues from \citet{Kaplan:20} and \citet{Brown:20short}.

Section~\ref{sec:ExpLanguage} expands the investigation to a larger-scale language modeling experiment. Here, xLSTM and the leading models from Section~\ref{sec:ExpComparison} are evaluated after training on 300B tokens from SlimPajama~\citep{Soboleva:23}. This section first assesses each model’s generalization to longer context lengths, then measures validation perplexity and downstream task performance as outlined by~\citep{Sutawika:24}. Subsequently, we test the models on 571 distinct text domains using the PALOMA language benchmark~\citep{Magnusson:23arxivshort}. Finally, the scaling characteristics of all methods are revisited in the context of substantially increased (20-fold) training data.

\label{sec:xlstm_blocks_notation}
Throughout all experimental setups, we denote xLSTM[$a$:$b$] to specify the proportion $a/b$ of mLSTM-based to sLSTM-based xLSTM blocks. To illustrate, xLSTM[7:1] indicates that among eight blocks, seven utilize mLSTM-based architecture while one employs an sLSTM-based design. In cases where the overall number of blocks is set to 48, this configuration corresponds to 42 mLSTM-based blocks and 6 sLSTM-based blocks. Additionally, in every experiment, we incorporate pre up-projection blocks for mLSTM and post up-projection blocks for sLSTM.

\subsection{Synthetic Tasks and Long Range Arena}
\label{sec:ExpSyntethic}
To begin, we evaluate how well xLSTM's novel exponential gating with memory mixing handles formal language tasks~\citep{Deletang:23}. Subsequently, we investigate the capabilities of xLSTM's new matrix memory mechanism on the Multi-Query Associative Recall benchmark~\citep{Arora:23arxiv}. Lastly, we analyze xLSTM's effectiveness in processing extended input sequences within the Long Range Arena framework~\citep{Tay:21}.

\paragraph{Evaluation of xLSTM's Exponential Gating and Memory Mixing Mechanisms.} We assess the effectiveness of xLSTM's newly introduced exponential gating combined with memory mixing, features designed to address state tracking tasks~\citep{Merrill:24, Merrill:23}. To do so, we adapt and build upon the formal language benchmarks from \citet{Deletang:23}, modifying them to support training across sequences of multiple lengths to facilitate extrapolation studies. Comprehensive details regarding each task, along with extended experimental results, are provided in Appendix~\ref{sec:appExpSynthetic-formalLang}. 

Our experimental setup benchmarks xLSTM against several architectures, including Transformers, State Space Models, and traditional Recurrent Neural Networks. We evaluate each method’s accuracy on task-specific tokens, reporting scores normalized from 0 (random guess) to 1 (flawless performance). Specifically, we compare 2-block implementations of various models: xLSTM[0:1] (representing sLSTM alone), xLSTM[1:0] (corresponding to mLSTM only), xLSTM[1:1], Llama, Mamba, RWKV, Retention, Hyena, LSTM, and LSTM incorporated within Transformer blocks (LSTM (Block)). The corresponding experimental results are presented in Figure~\ref{fig:formal-main}.

It is observed that models like Transformers or State Space Models lacking memory mixing and hence without state tracking abilities, are unable to solve formal language tasks such as the parity (regular grammar) problem. These findings are consistent with prior work indicating that, fundamentally, Transformers and State Space Models have inferior computational power relative to RNNs~\citep{Merrill:24, Merrill:23, Deletang:23}.

\paragraph{Test of xLSTM's Memory Capacities on Associative Recall Tasks.} This study evaluates the new matrix memory component of xLSTM with respect to its ability to retain information on the Multi-Query Associative Recall task~\citep{Arora:23arxiv}. In each trial, sequences are constructed using key-value pairs sampled randomly from an extensive vocabulary, requiring the model to store these associations for future retrieval. To increase task complexity, we scale the number of key-value pairs up to 256 and allow the context window to extend to 2048 tokens, thereby providing a more comprehensive challenge to assess the memory capabilities of various architectures. We conduct a comparative analysis involving 2-block versions of Llama, Mamba, RWKV-5, RWKV-6, xLSTM[1:1], and xLSTM[1:0], and evaluate each model’s accuracy in recalling the stored pairs. Given that Transformer-based models (such as Llama) are known for their memory scaling exponentially with coding dimensionality~\citep{Ramsauer:21}, they represent the benchmark for performance in this scenario. The outcomes are illustrated in Figure~\ref{fig:mqar-main}. 
Among non-Transformer baselines, xLSTM[1:1] achieves the highest recall accuracy, even at lower model sizes. Notably, the addition of the sLSTM block appears to enhance rather than impair memory performance, with pronounced benefits at the maximum challenge of recalling 256 key-value associations. Further findings are discussed in Appendix~\ref{sec:appExpSynthetic-mqar}, where extrapolation studies show that xLSTM’s expanded memory architecture supports generalization to context lengths surpassing those encountered during training.

\paragraph{Test of xLSTM's Long Context Capabilities on Long Range Arena.} To evaluate the ability of xLSTM to process extended sequences and broader contexts, we conduct a comparison of various approaches on the Long Range Arena~\citep{Tay:21}. Across all tasks, xLSTM consistently achieves robust results, indicating that its architecture is highly effective for diverse long context challenges.

Additional information can be found in Appendix~\ref{sec:appExpSynthetic-lra}.

\subsection{Method Comparison and Ablation Study}
\label{sec:ExpComparison}

In order to investigate our central research question—specifically, the capabilities of our newly proposed LSTM variants at larger scales in language modeling—we conduct experiments training xLSTMs, Transformers, State Space Models, along with other approaches, using 15B tokens sourced from SlimPajama, all within an auto-regressive framework. The resulting models are then evaluated on the validation set, and we carry out ablation analyses on the xLSTM architectures.

\paragraph{Comparing xLSTM to Other Methods.} To facilitate a fair evaluation, all models were trained on a dataset of 15B tokens sourced from SlimPajama~\citep{Soboleva:23}. Model performance is assessed using validation set perplexity. The following approaches are considered: our xLSTM, GPT-3 (Transformer)~\citep{Brown:20short}, Llama (Transformer)~\citep{Touvron:23arxiv}, H3 (SSM)~\citep{Fu:23}, Mamba (SSM)~\citep{Gu:24arxiv}, RWKV-4 (RNN)~\citep{Peng:23arxivshort}, RWKV-5 (RNN)~\citep{Peng:24arxivshort}, RWKV-6 (RNN)~\citep{Peng:24arxivshort}, GLA (linear Transformer)~\citep{Yang:23arxiv}, HGRN (RNN)~\citep{Qin:23}, HGRN2 (RNN)~\citep{Qin:24arxiv}, RetNet (linear Transformer)~\citep{Sun:23arxiv}, Hyena (linear Transformer)~\citep{Poli:23}, and xLSTM variants [1:0] and [7:1] (refer to Section~\ref{sec:xlstm_blocks_notation}). 

All models employ mixed precision training; however, for RWKV-5, RWKV-6, GLA, and HGRN2, the training did not leverage PyTorch’s automatic mixed precision routines (see Appendix Section~\ref{sec:appGenTrainProc} for more information). The compared architectures fall into three main classes: (a) Transformers, (b) State Space Models (SSMs), and (c) Recurrent Neural Networks (RNNs) alongside linear Transformers, the latter utilizing linearized mechanisms in place of traditional Transformer attention. The configurations align with a GPT-3 architecture containing 350M parameters, specifically using an embedding dimension of 1024 and 24 residual blocks. GPT-3 uniquely shares weights between the token and output embeddings, resulting in a lower parameter count than others. 

Empirical results presented in Table~\ref{tab:spaj15b_model_results} demonstrate that xLSTM yields superior validation perplexity relative to all baselines. Further analysis is available in Appendix~\ref{sec:appExpComparison}. Scaling trends depicted in Figure~\ref{fig:temp_sclaw15B} suggest that the advantages of xLSTM persist as model sizes increase.

\paragraph{Ablation Studies.}
As shown in Table~\ref{tab:spaj15b_model_results} and Figure~\ref{fig:temp_sclaw15B}, xLSTM delivers outstanding language modeling performance when trained on 15B SlimPajama tokens. To systematically analyze the modifications from the standard LSTM to xLSTM, we incrementally adapt a basic LSTM architecture toward xLSTM. The process begins with embedding LSTM layers into pre-LayerNorm residual frameworks. Next, we augment the architecture to include a post up-projection block. The final enhancement involves introducing exponential gating mechanisms and matrix memory.

Table~\ref{tab:ablstudies} (top) presents the findings.

The ablation studies indicate that the notable gains in performance can be ascribed to both the exponential gating mechanism and the incorporation of matrix memory. Furthermore, considering the crucial role of gating in both RNNs and State Space Models, we examine the effects of varying gating strategies. As summarized in Table~\ref{tab:ablstudies} (bottom), our results demonstrate that making each gate trainable and responsive to the input incrementally enhances model performance. Further analysis regarding specific backbone components can be found in Appendix~\ref{sec:appAblDetails}.

\subsection{xLSTM as Large Language Model}
\label{sec:ExpLanguage}

This investigation concludes with extensive language modeling experiments, where we assess xLSTM's capabilities as a large language model (LLM). To this end, we expand the training dataset, utilizing 300B tokens sourced from SlimPajama. This token count aligns with the amount used in prior works such as Mamba~\citep{Gu:24arxiv} and Griffin~\citep{De:24arxiv}.

We conduct a comparative study of xLSTM against RWKV-4, Llama, and Mamba—each representing a distinct architectural family as discussed in Section~\ref{sec:ExpComparison}. RWKV-4 serves as our baseline for RNNs, as suitable training precision settings for RWKV-5, RWKV-6, and HGRN2 were not determined until after the commencement of the 300B token training runs (refer to Appendix~\ref{sec:appGenTrainProc} for further details). We train models at various parameter counts (125M, 350M, 760M, and 1.3B) and evaluate all models in terms of their ability to extrapolate to longer sequences, as well as their accuracy on the validation set. In addition, we examine performance across several downstream tasks, measure their language modeling capabilities on the 571 domains included in the PALOMA benchmark, and analyze how each model family scales.

\paragraph{Sequence Length Extrapolation.}
\label{sec:spaj300B_ctx_extrapolation}
We begin by evaluating how well large-scale 1.3B parameter models—specifically, xLSTM, RWKV-4, Llama, and Mamba—can generalize to longer sequence lengths. These models are initially trained using a context window of 2048 tokens but are subsequently assessed with extended contexts up to 16384 tokens. The corresponding results are illustrated in Figure~\ref{fig:spaj300B_seq_extrapolation}. Notably, xLSTM models consistently achieve lower perplexity across the increased context lengths compared to the other architectures.

\paragraph{Validation Perplexity and Downstream Tasks.}
\label{sec:spaj_downstream_eval}
Next, across every model scale, we assess the capabilities of xLSTM, RWKV-4, Llama, and Mamba models using the SlimPajama validation dataset for next token prediction. Additionally, we examine their effectiveness on downstream benchmarks designed to evaluate common sense reasoning abilities.

The third column in Table~\ref{tab:lmeval} shows the perplexity values on the validation set for various approaches. xLSTM[1:0] and xLSTM[7:1] consistently achieve the lowest perplexities across all examined model sizes. The remaining columns in Table~\ref{tab:lmeval} report results on downstream evaluation benchmarks. In nearly all tasks and for every model size considered, xLSTM outperforms competing methods, with the sole exception occurring on the ARC task, where Mamba occasionally attains superior performance. Further information can be found in Appendix~\ref{sec:appExpLanguage}.

\paragraph{Performance on PALOMA Language Tasks.} Third, we evaluate the next-token prediction accuracy of xLSTM, RWKV-4, Llama, and Mamba models across all parameter scales using the PALOMA language suite~\citep{Magnusson:23arxivshort}. We report perplexity scores for next token prediction across 571 text domains, encompassing sources such as nytimes.com and the r/depression subreddit on Reddit. In Table~\ref{tab:ppleval}, token prediction perplexity is organized into broad language modeling (first seven columns) and detailed domain-specific benchmarks (last five columns). On these language benchmarks, xLSTM[1:0] consistently achieves lower perplexity compared to xLSTM[7:1].

In 568 out of 571 (99.5\%) text domains, xLSTM[1:0] achieves a lower perplexity than Mamba. Compared to Llama, it outperforms in 486 out of 571 (85.1\%) domains, and it yields a lower perplexity than RWKV-4 in 570 out of 571 (99.8\%) cases. Further information can be found in Appendix~\ref{sec:appExpLanguage}.

\paragraph{Scaling Laws.} Fourth, we investigate the scaling properties characterized by power-law trends, which are crucial for predicting how model performance changes with increasing model size~\citep{Kaplan:20,Brown:20short}. Figure~\ref{fig:temp_sclaw300B} illustrates these scaling patterns. While all examined models demonstrate comparable scaling tendencies, they are differentiated by constant performance shifts. Among them, RWKV-4 consistently shows the lowest performance, trailed by Llama and then Mamba. Notably, xLSTM outperforms Mamba by roughly the same relative amount that Mamba surpasses Llama. This scaling pattern suggests that, as model sizes expand, xLSTM is likely to maintain its advantage over both Transformer-based and State-Space models.

\textbf{Generation Times and Maximal Throughput.} In Figure~\ref{fig:inference_time_generation}, we report the measured text generation latency, while Figure~\ref{fig:inference_time_decoding_throughput} (left) presents the maximal throughput obtained with various xLSTM architectures at the 1.3B parameter scale. These results are contrasted against equivalently sized models of Mamba, Llama, and RWKV available on HuggingFace, and we include a static key-value cache for the Llama baseline. During experimentation, we were unable to perform full cache compilation for the Transformer Model or utilize \texttt{torch.compile} for the Mamba model, as these functionalities were not yet supported. All models in the text generation assessment use a batch size of 1 and a pre-fill of 16 tokens—settings that particularly favor Transformer inference.
As depicted in Figure~\ref{fig:inference_time_generation}, the inference time for xLSTM, Mamba, and RWKV-4 exhibits a linear relationship with sequence length, in contrast to the quadratic complexity observed in Llama. For decoding throughput, we evaluate across multiple batch sizes and pre-fill values specifically for Llama. 
In Figure~\ref{fig:inference_time_decoding_throughput} (right), the results illustrate that xLSTM accommodates substantially larger batch sizes than Llama due to its constant memory requirements, resulting in superior throughput performance.

\section{Limitations}

(i) Unlike mLSTM, the memory integration mechanism in sLSTM inherently precludes parallel operations, thus preventing efficient parallel implementation. Despite this, we engineered an accelerated CUDA kernel for sLSTM, which currently achieves performance at under twice the runtime of our parallel mLSTM variant. (ii) It is important to note that the existing CUDA kernels for mLSTM have not undergone thorough optimization, making them roughly four times slower than state-of-the-art solutions like FlashAttention or the scan operation employed in Mamba. Substantial speed improvements are plausible if the kernels are optimized similarly to FlashAttention. (iii) mLSTM's matrix-based memory imposes significant computational demands because it requires processing $d \times d$ matrices. Nonetheless, the processes of memory updating and retrieval remain parameter-free and benefit from efficient parallelization through conventional matrix operations, resulting in only a modest increase in actual computation time despite the added complexity. (iv) Careful selection of initial values for the forget gates is essential. (v) Since the size of the matrix memory does not depend on sequence length, extending sequence length for broader contexts may cause memory saturation in longer contexts. However, this effect has not proved problematic for sequences up to 16k tokens, as demonstrated in Section~\ref{sec:spaj300B_ctx_extrapolation}. (vi) Given the prohibitive computational cost associated with large-scale language model experiments, we have not conducted thorough architecture or hyperparameter tuning, particularly for the larger xLSTM configurations. We expect that achieving optimal xLSTM performance will require a rigorous and systematic optimization effort.

\section{Conclusion}
We have provided a partial response to our initial inquiry: To what extent can language modeling be advanced by increasing the scale of LSTM architectures to billions of parameters? Our findings suggest that LSTMs, when scaled in this way, can achieve performance levels comparable to those of present-day techniques such as Transformers and State Space Models. By introducing exponential gating with memory mixing and designing a novel memory mechanism, we have extended LSTM into xLSTM. In empirical comparisons on language modeling benchmarks, xLSTM models yield results that rival leading models, including both Transformers and State Space Models. Analyses of scaling laws reveal that expanding xLSTM to larger sizes positions it as a formidable alternative to today's Transformer-based Large Language Models. Beyond language modeling, xLSTM holds significant promise for influencing domains such as Reinforcement Learning, Time Series Prediction, and physical system modeling.

\end{document}